\section{Introduction}

Computer-use agents run a closed loop. They read a screenshot, act, and repeat until they declare the task done or run out of steps. A recurring way to fail is to declare it done when it is not: the agent announces success, the episode ends, and the environment's checker returns zero \citep{legacyworld2026}. Figure~\ref{fig:teaser}a shows one such episode. The agent writes that it cannot finish without a password, then reports that the installation ``has been initiated successfully''. Users of Holo report the same failure in deployment. In an open issue on H Company's desktop client, screen capture fails while a Windows elevation prompt is shown, 39 times in a row over 19 seconds; the agent keeps typing and clicking blind, then reports the task completed, while the author verified that nothing was done and an unsaved document was in the foreground \citep{holodesktop2026issue15}. In the author's words, ``The agent has no `I cannot see' state.'' % source: https://github.com/hcompai/holo-desktop-cli/issues/15 (read with gh, 2026-11-07); research_notes/Problèmes frontier de H/h_user_feedback.md item A.1


Interpretability offers a tempting repair. If a behaviour is mediated by one direction in the residual stream, projecting that direction out of every matrix that writes into the stream removes the behaviour, with no gradient step \citep{arditi2024refusal}. ``Weight-space programming'' recently used this to remove a looping habit from a 0.8B language model by editing 131{,}072 weights, and attributed the gain to removing an obstacle that hid an ability the model already had \citep{takehi2026wsp}. % source: BRIEF_neurochirurgie_holo_LOCAL.md section 0 (summary of the MIT post)
False success looks like the ideal target for this recipe. It is frequent (Fig.~\ref{fig:teaser}b) and it is a single discrete decision, the call to the \texttt{answer} tool, of the kind that single directions are known to steer \citep{wu2026toolcalling,chen2026toolrate}.

This paper is about the gap between finding such a direction and proving that it causes the behaviour. A probe that separates failing from healthy steps shows correlation; causation needs an intervention and a behavioural measurement \citep{belinkov2021probing,hewitt2019control}. In GUI agents that gap is wide, because the obvious signals are confounded with the position in the episode and because agent outcomes are noisy enough that an untreated re-run looks like an effect.

\paragraph{Contributions.} We report only what we measured.
\begin{enumerate}
\item \textbf{A diagnosis of a real agent on real desktops} (\S\ref{sec:diag}). Holo-3.1-4B solves \oswNSuccess/\oswNScored{} tasks of a frozen OSWorld split. Of \oswNFail{} failures, \oswCatFalseSuccess{} are false successes, \oswCatStepLimit{} hit the step limit and \oswCatLoop{} loop; \oswCatOther{} is unclassified. Re-sampled, the model says ``done'' on false and true completions alike (\S\ref{sec:think}).
\item \textbf{An exact replay pipeline} (\S\ref{sec:replay}) from a vLLM-served agent to per-layer activations in \texttt{transformers}, with message-level and token-level fidelity checks.
\item \textbf{An anti-confound protocol, motivated by three measured traps} (\S\ref{sec:traps}): position-confounded directions, regression to the mean between runs, and detection without causation. Each would have produced a false positive here.
\item \textbf{A causal record of every intervention we ran} (\S\ref{sec:causal-results}), as pre-registered pilot screens, with a layer-by-position causal map of the tool choice. None shows a large effect, and the screens could find only large ones: the gate flags a fall to \gateMaxAnswers/8 false answers or fewer, and one of the two false-success screening contexts is an infeasible task, so the evidence on feasible false successes rests on a single context and a few samples per arm (Appendix~\ref{app:infeasible}). The tool call follows the text the model has already written, and the edits that move it lower true completions too. On a strictly feasible cohort, a model merge lowers them as much as false completions, a lever without a discriminator, and the selectivity of a closed-form edit is not established either way. A detector of imminent completion beats position but does not separate false from true completions (\S\ref{sec:perceive}). No arm run online beats the run-to-run difference of the original (\S\ref{sec:online}). For the model's developer, the consequence is that a fix has to come from an external verifier or from training, not from weight surgery (\S\ref{sec:discussion}).
\end{enumerate}
Code, scripts and the provenance of every number: \url{https://github.com/martinbon39/ComputerUseHackaton} (\texttt{holo-neurochir/}). Raw outputs: released upon publication.

\placefig{fig-pipeline}
